\documentclass[10pt,a4paper]{amsart}
\usepackage{amsmath,amssymb,amsthm,graphicx}
\usepackage{verbatim}
\newenvironment{reference}{\comment}{\endcomment}
%\newenvironment{reference}{}{}
\newenvironment{slogan}{\comment}{\endcomment}
\newenvironment{history}{\comment}{\endcomment}

% For commutative diagrams we use Xy-pic
\usepackage[all]{xy}

% We use 2cell for 2-commutative diagrams.
\xyoption{2cell}
% Theorem environments.
%
\theoremstyle{plain}
\newtheorem{theorem}[subsection]{Theorem}
\newtheorem{proposition}[subsection]{Proposition}
\newtheorem{lemma}[subsection]{Lemma}

\theoremstyle{definition}
\newtheorem{definition}[subsection]{Definition}
\newtheorem{example}[subsection]{Example}
\newtheorem{exercise}[subsection]{Exercise}
\newtheorem{situation}[subsection]{Situation}

\theoremstyle{remark}
\newtheorem{remark}[subsection]{Remark}
\newtheorem{remarks}[subsection]{Remarks}

\numberwithin{equation}{subsection}

% Macros
%
\def\lim{\mathop{\mathrm{lim}}\nolimits}
\def\colim{\mathop{\mathrm{colim}}\nolimits}
\def\Spec{\mathop{\mathrm{Spec}}}
\def\Hom{\mathop{\mathrm{Hom}}\nolimits}
\def\Ext{\mathop{\mathrm{Ext}}\nolimits}
\def\SheafHom{\mathop{\mathcal{H}\!\mathit{om}}\nolimits}
\def\SheafExt{\mathop{\mathcal{E}\!\mathit{xt}}\nolimits}
\def\Sch{\mathit{Sch}}
\def\Mor{\mathop{\mathrm{Mor}}\nolimits}
\def\Ob{\mathop{\mathrm{Ob}}\nolimits}
\def\Sh{\mathop{\mathit{Sh}}\nolimits}
\def\NL{\mathop{N\!L}\nolimits}
\def\CH{\mathop{\mathrm{CH}}\nolimits}
\def\proetale{{pro\text{-}\acute{e}tale}}
\def\etale{{\acute{e}tale}}
\def\QCoh{\mathit{QCoh}}
\def\Ker{\mathop{\mathrm{Ker}}}
\def\Im{\mathop{\mathrm{Im}}}
\def\Coker{\mathop{\mathrm{Coker}}}
\def\Coim{\mathop{\mathrm{Coim}}}

% Boxtimes
%
\DeclareMathSymbol{\boxtimes}{\mathbin}{AMSa}{"02}

%
% Macros for moduli stacks/spaces
%
\def\QCohstack{\mathcal{QC}\!\mathit{oh}}
\def\Cohstack{\mathcal{C}\!\mathit{oh}}
\def\Spacesstack{\mathcal{S}\!\mathit{paces}}
\def\Quotfunctor{\mathrm{Quot}}
\def\Hilbfunctor{\mathrm{Hilb}}
\def\Curvesstack{\mathcal{C}\!\mathit{urves}}
\def\Polarizedstack{\mathcal{P}\!\mathit{olarized}}
\def\Complexesstack{\mathcal{C}\!\mathit{omplexes}}
% \Pic is the operator that assigns to X its picard group, usage \Pic(X)
% \Picardstack_{X/B} denotes the Picard stack of X over B
% \Picardfunctor_{X/B} denotes the Picard functor of X over B
\def\Pic{\mathop{\mathrm{Pic}}\nolimits}
\def\Picardstack{\mathcal{P}\!\mathit{ic}}
\def\Picardfunctor{\mathrm{Pic}}
\def\Deformationcategory{\mathcal{D}\!\mathit{ef}}
\usepackage{cleveref}
\crefname{originalsection}{original section}{original sections}
\Crefname{originalsection}{Original section}{Original sections}
\makeatletter
\newlabel{limits-section-limits}{{01YV}{source}{}{sourcepointer-1}{}}
\newlabel{limits-section-limits@cref}{{[originalsection][0][]01YV}{[1][0][]source}{}{}{}}
\makeatother
\crefname{originallemma}{original lemma}{original lemmas}
\Crefname{originallemma}{Original lemma}{Original lemmas}
\makeatletter
\newlabel{derived-lemma-two-ss-complex-functor}{{015J}{source}{}{sourcepointer-2}{}}
\newlabel{derived-lemma-two-ss-complex-functor@cref}{{[originallemma][0][]015J}{[1][0][]source}{}{}{}}
\makeatother
\title{Higher direct images commute with filtered colimits of bounded below complexes}
\author{The Stacks Project Authors}
\date{stacks-project:etale-cohomology.tex:local-snapshot-20261011; GNU FDL 1.2 or later}
\begin{document}
\maketitle
\noindent\textit{Editorial note.} All delivered mathematical text is the author's own, copied from the pinned Stacks Project chapter file etale-cohomology.tex; no statement or proof step was rewritten, reordered or removed.\par\medskip
\noindent\textit{Editorial note.} The proof is the author's complete argument: both spectral sequences are displayed and the comparison is justified by the bounded below hypothesis. The closing sentence 'Some details omitted' concerns the bookkeeping and is preserved verbatim in the counted span, together with the author's remark that stupid truncations give an alternative argument.\par\medskip
\noindent\textit{Editorial note.} Every labelled reference occurring inside the counted spans is carried in assembly.external\_references as a non-load-bearing navigation pointer (same-chapter ones at their definition site, cross-chapter ones at the invocation inside the counted span), so no load-bearing mathematics is replaced by a stub. No bibliography citation occurs inside the counted spans.\par\medskip
\noindent\textit{Editorial note.} The checked chapter list for same-name collisions is the set of stacks chapters staged in this repository: chow, cohomology, divisors, etale-cohomology, flat, more-algebra, more-morphisms, morphisms, spaces-more-morphisms, varieties together with this card's chapter etale-cohomology.\par\medskip
\noindent\textit{Editorial note.} The preamble is the chapter's own native preamble excerpt already used for this source: the Xy-pic setup (native preamble lines 14-18) and the native theorem declarations with the mathematical macros (native preamble lines 154-217).\par\medskip
\noindent\textit{Allowed standard prerequisites (editorial).}\begin{itemize}
\item The degreewise relative colimit result. Same chapter, Lemma lemma-relative-colimit-general. Supplies the interchange of filtered colimits with the higher direct images of the individual sheaves.
\item The spectral sequence of a complex under a functor. Derived Categories, Lemma derived-lemma-two-ss-complex-functor. Supplies the two spectral sequences whose first pages are compared.
\end{itemize}\medskip
\ifdefined\hypertarget\hypertarget{sourcepointer-1}{}\fi
\noindent\textit{Source reference.} Original section 01YV: sources/stacks-etale-cohomology/pinned-original-source.tex, source lines 7573--7573. Navigation pointer invoked at proof line 7573; the counted span contains the pointer only and the cited result is not reproduced inside it, so the pointer is registered as a non-load-bearing navigation stub in the established form for Stacks material. No bibliography citation is involved.\par\medskip
\ifdefined\hypertarget\hypertarget{sourcepointer-2}{}\fi
\noindent\textit{Source reference.} Original lemma 015J: sources/stacks-etale-cohomology/pinned-original-source.tex, source lines 7819--7819. Navigation pointer invoked at proof line 7819; the counted span contains the pointer only and the cited result is not reproduced inside it, so the pointer is registered as a non-load-bearing navigation stub in the established form for Stacks material. No bibliography citation is involved.\par\medskip
\begin{lemma}
\label{lemma-relative-colimit-general}
Let $I$ be a directed set. Let $g_i : X_i \to S_i$ be an inverse system of
morphisms of schemes over $I$. Assume $g_i$ is quasi-compact and
quasi-separated and for $i' \geq i$ the transition morphisms
$f_{i'i} : X_{i'} \to X_i$ and $h_{i'i} : S_{i'} \to S_i$ are affine.
Let $g : X \to S$ be the limit of the morphisms $g_i$, see
Limits, Section \ref{limits-section-limits}.
Denote $f_i : X \to X_i$ and $h_i : S \to S_i$ the projections.
Let $(\mathcal{F}_i, \varphi_{i'i})$ be a system of sheaves
on $(X_i, f_{i'i})$. Set $\mathcal{F} = \colim f_i^{-1}\mathcal{F}_i$. Then
$$
R^p g_* \mathcal{F} =
\colim_{i \in I} h_i^{-1}R^p g_{i, *} \mathcal{F}_i
$$
for all $p \geq 0$.
\end{lemma}

\begin{lemma}
\label{lemma-relative-colimit-general-really-complexes}
Let $I$, $g_i : X_i \to S_i$, $g : X \to S$, $f_{ii'}$, $f_i$, $g_i$, $h_i$
be as in Lemma \ref{lemma-relative-colimit-general}.
Let $\mathcal{F}_i^\bullet$ be a complex of abelian sheaves on
$X_{i, \etale}$. Let $\varphi_{i'i} : f_{i'i}^{-1}\mathcal{F}_i^\bullet \to
\mathcal{F}_{i'}^\bullet$ be a map of complexes on $X_{i, \etale}$
such that $\varphi_{i''i} = \varphi_{i''i'} \circ f_{i'' i'}^{-1}\varphi_{i'i}$
whenever $i'' \geq i' \geq i$. Assume there is an integer $a$ such that
$\mathcal{F}_i^n = 0$ for $n < a$ and all $i \in I$. Then
$$
R^pg_*(\colim f_i^{-1}\mathcal{F}_i^\bullet) =
\colim_{i \geq 0} h_i^{-1}R^pg_{i, *}\mathcal{F}_i^\bullet
$$
for all $p \in \mathbf{Z}$.
\end{lemma}

\begin{proof}
This is a consequence of Lemma \ref{lemma-relative-colimit-general}. Set
$\mathcal{F}^\bullet = \colim f_i^{-1}\mathcal{F}_i^\bullet$.
The lemma tells us that
$$
\colim_{i \in I} h_i^{-1}R^pg_{i, *}\mathcal{F}_i^n = R^pg_*\mathcal{F}^n
$$
for all $n, p \in \mathbf{Z}$. Let us use the spectral sequences
$$
E_{1, i}^{s, t} = R^tg_{i, *}\mathcal{F}_i^s \Rightarrow
R^{s + t}g_{i, *}\mathcal{F}_i^\bullet
$$
and
$$
E_1^{s, t} = R^tg_*\mathcal{F}^s \Rightarrow
R^{s + t}g_*\mathcal{F}^\bullet
$$
of Derived Categories, Lemma \ref{derived-lemma-two-ss-complex-functor}.
Since $\mathcal{F}_i^n = 0$ for $n < a$ (with $a$ independent of $i$)
we see that only a fixed finite number of terms $E_{1, i}^{s, t}$
(independent of $i$) and $E_1^{s, t}$ contribute and
$E_1^{s, t} = \colim E_{i, i}^{s, t}$.
This implies what we want. Some details omitted.
(There is an alternative argument using ``stupid'' truncations
of complexes which avoids using spectral sequences.)
\end{proof}

\end{document}
